%% ─────────────────────────────────────────────────────────────────────────
\subsection{Industrial and Technical Context}
\label{sec:intro_context}
%% ─────────────────────────────────────────────────────────────────────────

Supply chain planning has become a multi-dimensional assessment problem. An operator
deciding which of two late shipments to expedite is simultaneously weighing cost, deadline
pressure, supplier reliability and, increasingly, carbon footprint, while most monitoring
and planning tools still process one criterion at a
time~\autocite{pinedo2012scheduling,chakrabortty2020efficient}. That mismatch between the
dimensionality of the decision and the dimensionality of the tooling is precisely why
practitioners fall back on judgment-based prioritisation, and why the quality of that
judgment has become an engineering concern rather than a purely human one.

This thesis was produced during an R\&D placement in the \gls{din} of the ALTEN Group, which
conducts applied research where digital twin technology, logistics optimisation and
sustainable transport meet. Its Smart Green Supply Chain
programme~\autocite{benzakour2024history} pursues three aims at once. The lean aim attacks
waste, meaning cost, delay and over-stock, through pull-flow logic and real-time scheduling.
The green aim attacks the carbon footprint of each transport leg, working to the \gls{glec}
Framework v3.2~\autocite{smartfreight2023glec} and ISO~14083:2023~\autocite{iso14083}. The
predictive aim is the one this thesis sits inside: anticipating disruption through
probabilistic modelling, in the manner of a digital supply chain
twin~\autocite{ivanov2020digital}.

Behind all three is Industry~4.0~\autocite{kagermann2013industry}, and its promise that
cyber-physical systems, the \gls{iot} and digital twins can turn a reactive supply chain into
a self-organising one. The Physical Internet~\autocite{montreuil2011toward} supplies the
setting for the simulation environment used here, since it imagines logistics networks that
route shipments through shared infrastructure the way a network routes packets.

%% ─────────────────────────────────────────────────────────────────────────
\subsection{The DISCO Digital Twin and Prior Research at ALTEN-DIN}
\label{sec:intro_disco}
%% ─────────────────────────────────────────────────────────────────────────

DISCO, the Supply Chain Illustrator, is a serious game and digital twin developed by the
DIN to simulate a complete supply chain for a fictitious satellite product. It involves 22
component suppliers across structural assemblies, optical systems and embedded
electronics~\autocite{baxas2025architecture}. Players take the role of supply chain buyers
and decide on supplier selection, transport routes and factory allocation.
Figure~\ref{fig:disco_map} shows the supplier network they work with and the three
dimensions their decisions are scored on, and Table~\ref{tab:disco_stack} summarises the
technology stack underneath it.

\begin{figure}[htbp]
    \centering
    \includegraphics[width=0.82\linewidth]{Images/1.Introduction/disco_map.png}
    \caption{The DISCO serious game network}
    \label{fig:disco_map}
\end{figure}

Its scoring engine applies Min-Max normalisation over the values of the current game
session, illustrated here on the carbon dimension:

\begin{equation}
    \text{ScoreCO}_2 = 1 + 9 \times \frac{\text{realCO}_2 - \text{co}_{2,\min}}
                                         {\text{co}_{2,\max} - \text{co}_{2,\min}},
    \label{eq:score_co2}
\end{equation}

with analogous formulae for cost and risk, on a scale from 1 for the best achievable
performance to 10 for the worst.

\begin{table}[htbp]
\centering
\caption{DISCO technology stack}
\label{tab:disco_stack}
\begin{tabularx}{\linewidth}{lX}
\toprule
\textbf{Component} & \textbf{Technology / Value} \\
\midrule
Backend API    & Spring Boot + Java~21 \\
Database       & PostgreSQL (Azure) \\
Frontend       & Angular \\
Authentication & Player key / Azure AD SSO OAuth2 \\
Scores computed & CO\textsubscript{2}, Cost, Risk per supplier and itinerary, on $[1,10]$ \\
\bottomrule
\end{tabularx}
\end{table}

This work sits within a body of prior DIN research. \textcite{baxas2025architecture}
defines the DISCO software architecture including the \gls{tts} and \gls{ttr}
formalisation. \textcite{baulain2024internet} developed the Physical
Internet routing module, providing transit times through the HERE Maps API with Vincenty
distances and mandatory \gls{hgv} pause compliance, together with the Itinerary
Performance Index (\gls{ipg}) normalised on $[0,1]$.
\textcite{kabouche2025scoring} developed the supplier risk scorer combining delay
history, geographic vulnerability and financial fragility.
\textcite{cordova2024green} developed the per-route CO\textsubscript{2}
quantification module using the GLEC and VECTO methodology.
\textcite{mondelice2025tms} developed the Transport Management System interface
providing real-time transit-time extraction.

These works form the organisational and thematic context of the present research. The
system delivered here, named SupplyScore and described in Section~\ref{sec:methods_impl},
is a standalone Python application that runs alongside these modules rather than inside
them. Their coupling is a stated perspective of the programme, not a delivered capability,
and this thesis makes no claim to the contrary.

%% ─────────────────────────────────────────────────────────────────────────
\subsection{Problem Statement}
\label{sec:intro_problem}
%% ─────────────────────────────────────────────────────────────────────────

This thesis addresses the layer that comes \emph{before} any planning or scheduling
decision, and which the classical literature has not fully investigated: the urgency
signal itself is systematically biased at the source, and must be measured and corrected
before it can safely inform any downstream
decision~\autocite{zhu2018mere,rusou2020psychology}.

The consequences are worse where capacity is shared. In an interconnected network, a biased
urgency signal does not stay local: it propagates, and misallocation grows as it travels.
Over-triage, declaring more urgency than the situation warrants, takes capacity from
whoever needed it more. Under-triage does the opposite and is more dangerous, because a risk
nobody has declared is a risk nobody is
watching~\autocite{li2020exploring,ivanov2020digital}.

Digital twin architectures and real-time decision support systems have demonstrated strong
capabilities in operational monitoring, disruption detection, resource reallocation and
order synchronisation~\autocite{ding2023realtime,leung2022digital,hrusovsky2021realtime,karam2022realtime,zhu2025industry5}.
They do so primarily by computing the physical operational state of the system. To the
best of the review presented in Section~\ref{sec:Literature}, existing systems do not
explicitly model the gap between a declared urgency and a computed operational urgency.

Figure~\ref{fig:urgency_gap} shows what that gap looks like over the life of a task. Two
curves are plotted against time. The step line is declared urgency, which moves only when
the operator submits a review, and the smooth curve is real urgency, which moves as the
operational state does. Reading left to right, the figure passes through four phases: the
two signals agree, then a first adequacy defect opens, then a second defect of the opposite
sign, then agreement returns. Shading tells the two defects apart. Where the step line sits
above the curve the operator is claiming more urgency than the state justifies, and the
enclosed area is false urgency. Where the curve sits above the step line the state is worse
than anyone has declared, and the enclosed area is hidden risk. Leader lines pointing at a
curve label that signal; leader lines pointing into a shaded area label the gap between
them.

\begin{figure}[htbp]
    \centering
    \includegraphics[width=0.82\linewidth]{Images/2.Problem/urgency_gap.png}
    \caption{Declared and real urgency, and the two mismatch pathologies}
    \label{fig:urgency_gap}
\end{figure}

Formally, let $U_d \in [0,1]$ denote the declared urgency and $U_r(t) \in [0,1]$ the real
urgency computed from operational data in the digital twin. The mismatch between the two
signals is

\begin{equation}
    \Delta U(t) = U_d - U_r(t) \in [-1,\, 1].
    \label{eq:delta_u}
\end{equation}

A tolerance band $\varepsilon \in [0,1]$ is introduced so that small deviations are not
treated as operationally meaningful. Over-declaration then corresponds to
$\Delta U(t) > \varepsilon$, under-declaration to $\Delta U(t) < -\varepsilon$, and
approximate agreement to $|\Delta U(t)| \le \varepsilon$. Table~\ref{tab:pathologies}
gives the operational reading of each condition.

\begin{table}[htbp]
\centering
\caption{Urgency pathologies and their consequences}
\label{tab:pathologies}
\small
\begin{tabularx}{\linewidth}{l >{\raggedright\arraybackslash}X >{\raggedright\arraybackslash}X}
\toprule
\textbf{Condition} & \textbf{Interpretation} & \textbf{Potential consequence} \\
\midrule
$\Delta U(t) > \varepsilon$ & Over-declaration of urgency, or false urgency & Excess
expediting, avoidable cost, possible delay for other flows \\[4pt]
$\Delta U(t) < -\varepsilon$ & Under-declaration of urgency, or hidden risk & Hidden
lateness or stockout risk, delayed intervention, downstream disruption \\[4pt]
$|\Delta U(t)| \le \varepsilon$ & Acceptable alignment & No material urgency mismatch \\
\bottomrule
\end{tabularx}
\end{table}

The two conditions need not cost the same. Weighting them asymmetrically, with
$\lambda_{\text{under}} \ge \lambda_{\text{over}}$, prices under-prioritisation above
over-prioritisation. Two bodies of work support doing so: operations research showing that
decision errors carry unequal downstream
consequences~\autocite{kim2020admission,diwa2020task}, and scheduling work that models
differentiated tardiness penalties instead of a uniform delay
cost~\autocite{baker1982dynamic,chakrabortty2022}. How much asymmetry is right is a
calibration question, answerable only in a specific setting.

%% ─────────────────────────────────────────────────────────────────────────
\subsection{Aim, Objectives and Research Questions}
\label{sec:intro_aim}
%% ─────────────────────────────────────────────────────────────────────────

The aim of this research is to establish whether the divergence between declared and
computed urgency is measurable to a standard that would justify acting on it.

That wording is deliberate on two counts. It commits to a falsifiable target rather than to
building a system: a score that cannot beat a trivial baseline against recorded outcomes has
not met the standard, however well engineered the system producing it. And it makes
measurability the question rather than assuming it, because the two signals being compared
are not independent by construction, and an architecture that lets the same person supply
both sides is measuring its own input.

Three objectives follow from that aim:

\begin{enumerate}[label=\textbf{O\arabic*.}]
    \item Define a multi-block, DAG-propagated real urgency $U_r(t)$ and a cognitive
          adequacy score $A(t) \in [0, 100]$ measuring the gap between the planner's
          declared urgency $U_d$ and the digital twin's operational state.
    \item Develop and implement a network propagation and criticality-driven
          prioritisation layer, combining PROMETHEE~II outranking with systematic shock
          simulation, to identify structurally critical nodes.
    \item Validate the resulting signal empirically, against recorded outcomes, and
          measure what happens to human declaration behaviour when the system's own
          forecast is placed in front of the declarant.
\end{enumerate}

Four research questions structure the scientific approach:

\begin{enumerate}[label=\textbf{RQ\arabic*}., leftmargin=*]
  \item \emph{Multi-block real urgency.} How can a physically and operationally grounded
  real urgency $U_r(t)$ be computed dynamically in a multi-tenant digital twin from
  heterogeneous KPI blocks, covering time, capacity, performance, risk, cost and
  environmental dimensions, and propagated coherently across a multi-tier supply chain
  graph?

  \item \emph{Cognitive adequacy formulation.} How should an adequacy score $A(t)$ be
  formulated to measure the alignment between declared urgency $U_d$ and computed real
  urgency $U_r(t)$ while asymmetrically penalising under-triage against over-triage, and
  does that score demonstrate predictive validity against recorded operational outcomes?

  \item \emph{Network prioritisation and criticality.} How can pairwise elicitation
  methods, namely AHP and the Fuzzy Best-Worst Method, be combined with outranking
  aggregation through PROMETHEE~II and systematic shock simulation to identify which nodes
  of a supply network are structurally critical?

  \item \emph{Operational validation and the effect of showing the forecast.} How can
  this framework be operationalised and empirically validated within a persistent digital
  twin exercised as a weekly-cadence serious game, and what happens to declared urgency
  once the declarant is shown the system's own prediction?
\end{enumerate}

The second half of RQ4 is not an afterthought. The campaign produced a result about the
declarant rather than only about the model, and reporting the model's scores without it
would be misleading (Section~\ref{sec:res_armB}).

%% ─────────────────────────────────────────────────────────────────────────
\subsection{Research Hypotheses}
\label{sec:hypotheses}
%% ─────────────────────────────────────────────────────────────────────────

These six hypotheses were pre-registered. They were written, with their numeric thresholds,
and frozen before the campaign of Section~\ref{sec:methods_campaign} produced any data, which
is why they appear here rather than next to the results they decide. They are reproduced
unchanged. Where the campaign later showed a threshold to be badly chosen, or a hypothesis to
be untestable on the scenario as scripted, that is reported as a finding in
Section~\ref{sec:hypothesis_status} and not repaired by rewriting the hypothesis.

Each states an observable outcome with a threshold the delivered system either meets or does
not, and each serves one of the research questions above.

\begin{description}
    \item[H1 (Cognitive latency):] Declared urgency follows computed real urgency with a
    delay of at least one weekly cycle on the majority of nodes. Test: the lagged
    correlation between $U_d$ and $U_r$ peaks at a lag of one week or more for at least
    five of the eight campaign nodes.
    \item[H2 (Early detection of hidden risk):] The hidden-risk signal
    $H = \max(U_r - U_d, 0)$ exceeds $0.3$ at the crisis-origin node several weeks before
    the rupture becomes publicly visible.
    \item[H3 (Structural criticality):] The systematic shock simulation ranks the true
    crisis origin first, by impact on the final customer, on at least 80\% of the
    informative, non-saturated weekly rounds.
    \item[H4 (Predictive validity):] At threshold $H \ge 0.5$ with a four-week horizon,
    the hidden-risk signal predicts recorded negative outcomes, meaning a missed
    milestone, a critical event or a node abandonment, with precision and recall both
    above $0.5$.
    \item[H5 (Urgency contagion):] On weeks when crisis news becomes public, operators of
    nodes that are not physically affected still raise their declared urgency,
    reproducing the mere-urgency effect~\autocite{zhu2018mere} at network scale.
    \item[H6 (Perception hysteresis):] Once the physical situation improves, declared
    urgency decays at less than half the rate of real urgency.
\end{description}

Table~\ref{tab:rq_hypothesis_map} maps each hypothesis to the question it serves and states
what its failure would mean, so that a negative result carries a consequence rather than
just a verdict.

\begin{table}[htbp]
\centering
\caption{Each hypothesis, its research question, and the cost of failure}
\label{tab:rq_hypothesis_map}
\small
\begin{tabularx}{\linewidth}{l >{\centering\arraybackslash}p{1.4cm} X}
\toprule
\textbf{Hypothesis} & \textbf{Serves} & \textbf{If it fails} \\
\midrule
H1 cognitive latency & RQ2 & Declared urgency tracks the physical signal without delay, so
there is no lag to correct and the adequacy layer has nothing to add. \\
H2 early detection & RQ2 & The hidden-risk signal arrives no earlier than the rupture it is
meant to anticipate, which makes it a description rather than a warning. \\
H3 structural criticality & RQ3 & The network ranking cannot be trusted to identify the
bottleneck, so criticality-driven prioritisation is unsafe to act on. \\
H4 predictive validity & RQ2 & The signal does not predict recorded outcomes, and the score
is a diagnostic curiosity rather than a decision aid. This is the hypothesis the aim of
Section~\ref{sec:intro_aim} stands or falls on. \\
H5 urgency contagion & RQ1, RQ4 & Operators do not import urgency from context, so
declaration bias is individual rather than systemic and does not propagate. \\
H6 perception hysteresis & RQ2 & Over-declaration corrects itself once the situation eases,
so the asymmetric penalty is mis-specified and $\lambda_{\text{under}}$ needs
recalibrating. \\
\bottomrule
\end{tabularx}
\end{table}

One omission from this register is worth flagging in advance, because
Section~\ref{sec:res_target_trap} measures its cost. The thresholds were pre-registered but
the definition of an adverse outcome was not, and that definition turns out to move the
measured result by more than the thresholds do.

Distinct from these testable hypotheses, the framework rests on stated modelling
assumptions, each verifiable in the delivered code: the network is a directed acyclic
graph, so there are no feedback loops and a linear topological sort is valid; reviews
follow a strict ISO weekly cycle; declared urgency is damped as it descends toward
suppliers while real risk is amplified as it ascends toward the customer, each scaled by a
per-arc coefficient; the six KPI blocks are treated as independent so that a probabilistic
aggregation is admissible; task statuses drive explicit override rules; and the cognitive
loss is deliberately asymmetric, with under-declaration priced at roughly twice the cost
of panic. These are design commitments, not empirical claims, and the limitations they
induce are catalogued in Section~\ref{sec:limitations}.

%% ─────────────────────────────────────────────────────────────────────────
\subsection{Scope and Delimitations}
\label{sec:intro_scope}
%% ─────────────────────────────────────────────────────────────────────────

Four boundaries are set explicitly, because each of them shapes what the results in
Section~\ref{sec:Results} can and cannot support.

\paragraph{The declarants are language-model agents, not human consultants.} The campaign
reported here was run with eight isolated language-model agents, one per role, each
confined to its own role card and turn briefings. This is an experimental design choice
with real consequences for external validity, discussed in
Sections~\ref{sec:methods_arms} and~\ref{sec:res_discussion}. The human campaign for which
the protocol was originally frozen has not been run.

\paragraph{The validation is a single scenario.} All results come from one pre-registered
scenario, a semiconductor shortage propagating through an eight-node satellite supply
chain over nineteen weekly turns. Statements about how the framework behaves are
statements about its behaviour on that scenario.

\paragraph{SupplyScore is standalone.} It does not integrate with DISCO, with an ERP, or with
the routing and carbon modules of Section~\ref{sec:intro_disco}. Every KPI reaching the
model arrives either through a manual weekly review or through a scripted campaign event.

\paragraph{The satellite observation work is an opening, not a second contribution.}
Section~\ref{sec:res_volume} reports a separate system that estimates warehouse capacity
from public satellite imagery. It is included because it addresses the freshness
limitation that the rest of this thesis documents about itself, and it is reported at the
depth that claim requires, without any assertion that the two systems have been connected.

%% ─────────────────────────────────────────────────────────────────────────
\subsection{Thesis Structure}
\label{sec:intro_structure}
%% ─────────────────────────────────────────────────────────────────────────

Four sections follow. Section~\ref{sec:Literature} reviews the streams that bear on the
problem and states the gap each one leaves. Section~\ref{sec:Methods} keeps three things
apart that are easily confused: the mathematical model, the choices made in building it into
software, and the experiment designed to test it. It closes with the method's own limits,
its reproducibility provisions and its ethics. Section~\ref{sec:Results} reports what was
measured, then diagnoses it. Section~\ref{sec:Conclusion} states what the work contributes,
what each hypothesis did, where the work stops, and what it opens.
